\documentclass[11pt]{article}
\usepackage[T1]{fontenc}
\usepackage[utf8]{inputenc}
\usepackage[a4paper,margin=2.4cm]{geometry}
\usepackage{amsmath,amssymb}
\pagestyle{plain}
\setlength{\parskip}{4pt}
\newcommand{\su}{\mathfrak{su}}

\begin{document}

\begin{center}
{\Large\bfseries Archive Erratum and Full Retraction (v3)}\\[5pt]
{\large Interface Lemmas for the Multiscale Proof of the\\
Lattice Yang--Mills Mass Gap}\\[6pt]
Lluis Eriksson \quad\textbullet\quad ai.viXra:2602.0052 \quad\textbullet\quad July 2026
\end{center}

\noindent\rule{\textwidth}{0.8pt}

\vspace{2pt}
\noindent
This version does two things. \textbf{First it restores the identity of the record}
(\S1), which since 2026-07-06 has displayed one paper's metadata while serving a
second paper's PDF, leaving this paper's title absent from the public index.
\textbf{Second, having restored it, it retracts the paper} (\S2--\S5). Not only its
terminal claim: an audit of every numbered statement found that \emph{none of the
three interface lemmas is established as printed}, that one auxiliary \emph{construction
claim inside a proof} is outright false, and that a second auxiliary proof is invalid
over its claimed parameter range. The manuscript follows unmodified so that the claims
and their refutations can be read against one another.

Throughout, a distinction is kept that this note's own earlier drafts twice failed to
keep: \emph{refuting a printed proof is not refuting the statement it was meant to
prove}. Where only the proof fails, that is what is said.

\section{Archive erratum: what this record has been serving}

\begin{center}
\begin{tabular}{|l|l|}
\hline
\textbf{Field} & \textbf{Actual provenance}\\
\hline
Record identifier & \texttt{2602.0052}\\
\hline
Title, abstract, category shown & \texttt{2602.0036v2}, \emph{Geodesic Convexity\ldots}\\
\hline
PDF served as v2 & \texttt{2602.0035v2}, \emph{Morse--Bott Spectral Reduction\ldots}\\
\hline
Legitimate content of the record & \texttt{2602.0052v1}, \emph{Interface Lemmas\ldots}\\
\hline
\end{tabular}
\end{center}

\noindent
Filiation was established from content, not filenames: the file served as v2 carries
the cover, subject, theorem numbering (3.1, 4.1) and explicit self-identification
(``\emph{version 2 records that\ldots}'') of the Morse--Bott manuscript, and is not
byte-identical to \texttt{2602.0036v2}, itself a genuine separate revision of its own
record. Only the metadata came from \texttt{2602.0036}; the file came from
\texttt{2602.0035}. A replacement creates a new version rather than overwriting the
old, so \emph{Interface Lemmas}, dated February 2026 and submitted here on
2026-02-12, is the authentic content of this record, not an orphan.

Between 2026-07-06 and the present correction, therefore, the title of this paper did
not resolve in the public index.

\section{The paper is the last document of its lane still claiming an unconditional gap}

The campaign that overwrote this record in July 2026 was, at that moment, retracting
the ``unconditional'' thesis across this lane. Because the write landed on the wrong
record, this paper never received its correction.

\texttt{2602.0053} v2 is retitled from ``\ldots and \emph{Unconditional} Mass
Gap\ldots'' to ``a \emph{Conditional and Windowed} Reduction''. \texttt{2602.0054} v2
states that ``\emph{the inter-block assembly requires the block Dobrushin condition
--- this is where {\rm(H-DOB-blk)} enters}'' and withdraws by name the claim that the
route ``\emph{bypasses any explicit Dobrushin contraction estimate}''.
\texttt{2602.0055} v2 lists (H-P0), (H-YGZ), (H-DOB-blk) and a volume window,
records that ``\emph{None of these is proved in the series}'', and declares ``\emph{the
retagging of the chain is complete}''. It was not: this record was outside it.

Accordingly \textbf{Corollary 7.3} (``Unconditional mass gap'',
$\Delta_{\rm phys}\ge m(\beta,N_c,d)>0$ uniformly in $L$), \textbf{Remark 6.4}
(``Dobrushin not needed''), and the closing sentence of the abstract are
\textbf{withdrawn}. The dating is independently checkable: the bibliography cites
[13], [15] and [17] under their pre-July titles, all three since retitled.

\section{Two invalid proofs, one of them resting on a false construction}

\subsection{Lemma A.1: the printed logarithm leaves the Lie algebra}

Lemma A.1 constructs $\log U:=V\operatorname{diag}(i\theta_1,\dots,i\theta_{N_c})V^{-1}$
with $\theta_j\in(-\pi,\pi]$, and Remark A.2 emphasises that this ``defines $\log U$
for every $U\in SU(N_c)$'' with ``no exceptional set discarded''. Take $U=-I\in SU(2)$,
whose eigenvalue arguments are both $\pi$:
\[
\log(-I)=i\pi I,\qquad \operatorname{tr}\log(-I)=2\pi i\neq0,
\qquad\text{so}\quad \log(-I)\notin\su(2).
\]
The output leaves the Lie algebra. This is not repaired by shifting one angle by
$2\pi$. The block map $Q_k=\exp\circ\operatorname{avg}\circ\log\circ\operatorname{mult}$
needs a logarithm equivariant under conjugation, because Lemma 2.2 asserts
$Q_k(g\cdot U)=g_{k+1}\cdot Q_k(U)$. But $-I$ is central: $g(-I)g^{-1}=-I$ for all $g$,
so equivariance would force $\log(-I)$ to commute with all of $SU(2)$, hence to be a
traceless multiple of the identity, hence $0$ --- and $\exp 0=I\neq-I$.

\textbf{No global, $\su(N)$-valued, conjugation-equivariant section of $\exp$ exists at
non-trivial central elements.}

\emph{What this does and does not refute.} It refutes the construction: the printed
map does not take values in $\su(N)$, and no repair of it can be simultaneously
$\su(N)$-valued, a section of $\exp$, and conjugation-equivariant. It therefore
refutes the printed justification that $Q_k$ is a globally defined, $SU(N)$-valued,
gauge-covariant map, and with it the printed proof of Lemma 2.2. It does \emph{not}
refute the statement of Lemma A.1, which asserts only that
$Q_k:\mathcal A_k\to\mathcal A_{k+1}$ is Borel measurable. A Borel $\su(N)$-valued
logarithm can be selected after branch corrections --- dropping equivariance, which
measurability does not require --- so the measurability claim is plausibly recoverable.
\textbf{Lemma A.1 is therefore not established; its construction is false.}

Upgrading this to a refutation of the statement would require exhibiting a fine
configuration on which the printed formula leaves $SU(N)$. The shape is visible: if a
block average runs over $M>1$ terms and one holonomy equals $-I$ while the others equal
$I$, the principal selection averages to $(i\pi/M)I$, whose exponential has determinant
$e^{2\pi i/M}\neq1$. Whether that holonomy pattern is realisable for the specific paths
$\Gamma_{c,x}$ of the block map is a question about its geometry, and is \emph{not}
settled here; the claim is not made.

Since Lemma 3.4 conditions on iterates of $Q_k$, the failure of the global covariance
justification reaches the Horizon Transfer chain as well.

\subsection{Lemma 6.2: the Holley--Stroock bound is applied in the wrong direction}

Lemma 6.2 asserts $\alpha_{\rm blk}\ge\frac{N_c}{4}e^{-C_{\rm fib}(\beta_0)}$ with
$C_{\rm fib}(\beta_0)=2\beta_0 n_{\rm plaq}+C_{\rm poly}$, proved from
$\operatorname{osc}(W_{k,B})\le2\beta\,n_{\rm plaq}+C_{\rm poly}$ by ``\emph{bounding
the oscillation using $\beta=\beta_0$}''. Substituting the smallest $\beta$ into an
upper bound that increases with $\beta$ does not preserve it. \textbf{The printed
proof is therefore invalid over the stated range} $\beta\ge\beta_0$.

\emph{This does not refute the statement.} Holley--Stroock supplies one sufficient
lower bound on the log-Sobolev constant; that this particular bound degenerates as
$\beta$ grows says nothing about the optimal constant, which some other argument might
still bound uniformly. Nor do three sampled values --- $\operatorname{osc}
(\beta=1,10,100)=3.0,\,30.3,\,303.1$, measured in \texttt{2602.0054} v2 --- prove that
$\inf_{\beta\ge\beta_0}e^{-\operatorname{osc}(\beta)}=0$; and the paper supplies only
an \emph{upper} bound on $\operatorname{osc}$, which cannot establish that
$\operatorname{osc}$ grows at all. The correct status is \textbf{not established over
the unbounded range; the same route does give a uniform bound on a bounded window}
$\beta_0\le\beta\le\beta_1$, with the constant depending on $\beta_1$.

\textbf{Windowing does not repair Lemma 6.3}, which additionally needs
$\delta_k<\alpha_{\rm blk}/4$ uniformly before applying
Yoshida--Gian\-nulis--Zegarlinski --- and that inter-block smallness is precisely the
content of (H-DOB-blk)/(H-YGZ), unproved in the series. An earlier draft of this note
said ``Lemma 6.3 stands only in the windowed form''; that was wrong, and is withdrawn.

\section{The three interface lemmas are not established as printed}

\subsection{Lemma C (fiber LSI with boundary)} Lemma 6.3 loses its per-block input
(\S3.2) and carries the inter-block hypothesis. \textbf{Conditional}, not established.

\subsection{Lemma A (Horizon Transfer, Lemma 3.11)} Two bridges are missing.

\emph{(i) The RCP/RG identification is asserted, not derived.} Lemma 3.4 states that the
regular conditional probability on a fibre of $Q^{(k)}$ has density proportional to
$e^{-S_k^{\rm eff}}$ against an ``induced Haar measure on the fibre'', and invokes
uniqueness of the RCP. A fibre of a nonlinear $Q^{(k)}$ is not a group and carries no
canonical Haar measure; a disintegration over it may involve Jacobians, gauge-fixing
coordinates and normalisations; and conditioning on $\pi\circ Q^{(k)}$ fixes a gauge
\emph{orbit}, not a representative. Uniqueness identifies two kernels only after both
are shown to disintegrate the same measure, which is what would have to be proved.

\emph{(ii) The horizon component is random but used as if fixed.} $Y_B$ (Definition 3.2)
is the connected large-field component containing $B$, and depends on the very fast
variables being integrated; the paper does not establish its measurability. Yet
Lemma 3.9 writes $\mu_k(L_k(B)\mid\mathcal G_{k+1})=T'_k(Y_B,(\bar U,0))F_B$, whose
left-hand side is $\mathcal G_{k+1}$-measurable and whose right-hand side is not. The
expected form is a sum over branches,
$\mu_k(L_k(B)\mid\mathcal G_{k+1})=\sum_{Y\ni B}T'_k(Y,(\bar U,0))F_{B,Y}$, requiring
measurability of each branch, disjointness, conditional normalisation, and
summability over all $Y\ni B$. The same confusion appears in Proposition 3.10, whose
case ``$Y_B=\emptyset$'' concludes that the conditional probability is $0$ --- but
whether $Y_B$ is empty depends on the fast variables, so conditioning on the slow field
cannot make it deterministically zero.

Moreover the per-component bound $T_k(Y)\mathbf 1\le e^{-cp_0(g_k)}$ does not control a
sum over all animals $Y\ni B$ without a size- or distance-dependent estimate and an
animal-counting argument. Finally, the statement of Lemma 3.11 declares $c$ to depend
only on $(d,N_c,L_{\rm RG},\gamma_0)$, while its proof takes $c=2/(1+\beta_0)$ from
input (viii); $\beta_0$ is missing from the dependency list.

\textbf{Not established.}

\subsection{Lemma B (uniform analyticity, Lemma 4.1)} The printed proof argues that
``$R^{(k)}$ and $B^{(k)}$ are summands of $S_k$ with disjoint supports; each inherits
the analyticity domain''. Analyticity of a sum does not give analyticity of the
summands; the sentence that follows, citing a direct construction of $B^{(k)}$, is the
argument that would have to be made, and is not carried out here. \textbf{Pending
literal verification against the cited source.}

\begin{sloppypar}
Its corollaries are worse. Input (v) reads
$|B^{(k)}(X)|\le O(1)\sum_{j\le k}|\Gamma^0_j\cap X|$ --- with \emph{no} spatial decay
factor, unlike input (iv) for $R^{(k)}$. Yet Corollary 4.2 concludes
$|\nabla_vB^{(k)}(X)|\le(C_Bk|X|/\hat\alpha_1)\,e^{-\kappa d_k(X)}$. A Cauchy estimate
converts a sup bound into a derivative bound; it cannot manufacture decay absent from
the input. Corollary 4.3 then sums over polymers using exactly that decay.
\textbf{Corollaries 4.2--4.3 are not established.}
\end{sloppypar}

\section{A further invalid inference: Lemma 5.6}

The proof of Lemma 5.6 argues that since $p_0(g)\ge c_0|\log g|^{1+\epsilon_0}$ and
$g_k\to0$, ``$p_0(g_k)\to\infty$ faster than any multiple of $k$'', hence
$(1+k)^2e^{-cp_0(g_k)}$ ``decays faster than any geometric sequence''. \textbf{The
displayed assumptions do not imply this.} The input on $p_0$ is a \emph{lower} bound,
so nothing in it prevents $p_0$ from being much larger --- but nothing in it forces
$p_0$ to be larger either, and the conclusion needs the latter.

A model compatible with every printed inequality makes this concrete. Take
\[
\beta_k=\beta_0+Ck,\qquad g_k=(\beta_0+Ck)^{-1/2},\qquad
p_0(g)=c_0|\log g|^{1+\epsilon_0} .
\]
This satisfies the running-coupling bound $\beta_k\le\beta_0+Ck$, satisfies
$g_k\to0$, and satisfies $p_0(g)\ge c_0|\log g|^{1+\epsilon_0}$ with equality. In it
\[
p_0(g_k)=O\bigl((\log k)^{1+\epsilon_0}\bigr),
\]
so $e^{-cp_0(g_k)}$ is subexponential --- superpolynomially small, but not geometric
in $k$ --- and the required bound $\le C_{\rm SF}^2L_{\rm RG}^{-(d-1)k}$ fails. The
geometric estimate therefore does not follow from the printed inputs; an additional
hypothesis, such as $p_0(g_k)\ge ck$ or the lane's explicit (H-P0), is required.
Lemma 5.7 loses its geometric input accordingly.

\section{Status}

\begin{center}
\begin{tabular}{|l|l|}
\hline
Archive filiation erratum (\S1) & Stands\\
\hline
Corollary 7.3, Remark 6.4, closing sentence & \textbf{Withdrawn}\\
\hline
Principal-angle $\log:SU(N)\to\su(N)$
 & \textbf{FALSE} as a construction (\S3.1)\\
\hline
Definition 2.1 / global block map $Q_k$
 & \textbf{Not established} as an $SU(N)$-valued\\
 & gauge-covariant map with the printed\\
 & logarithm (\S3.1)\\
\hline
Lemma A.1 (Borel measurability of $Q_k$)
 & Proof invalid, statement \textbf{not established};\\
 & plausibly recoverable without equivariance\\
\hline
Lemma 2.2 (gauge covariance of $Q_k$) & Not established globally (\S3.1)\\
\hline
Lemma 6.2 (per-block LSI)
 & Printed proof \textbf{invalid}; not established\\
 & over $\beta\ge\beta_0$; this route does give a\\
 & uniform bound on a bounded window (\S3.2)\\
\hline
Lemma 3.4, 3.9, Prop.\ 3.10, Lemma 3.11 & \textbf{Not established} (\S4.2)\\
\hline
Lemma 4.1 & Pending verification against the source\\
\hline
Corollaries 4.2--4.3 & \textbf{Not established} (\S4.3)\\
\hline
Lemma 5.6 & Invalid inference; needs (H-P0)\\
\hline
Lemma 5.7, Lemma 6.3, Theorem 7.1 & Conditional / not established\\
\hline
\end{tabular}
\end{center}

\noindent
\textbf{No numbered statement of this paper is claimed here to be false.} What is
claimed false is one construction inside a proof: the principal-angle map is not an
$\su(N)$-valued logarithm, and no equivariant repair of it exists. Everything else
above is a failure of a printed argument. The three interface lemmas may yet hold; so
may Lemma 6.2 over the unbounded range; and Borel measurability of a suitably
reconstructed block map is not claimed to be false. What this note records is that the
printed arguments do not establish them, and that the closure built on top of them does
not stand.

\vspace{4pt}
\noindent\rule{\textwidth}{0.8pt}

\noindent\footnotesize
\textbf{Provenance.} The mis-filing (\S1) was measured on 2026-07-29 by comparing every
record's archive metadata against the cover page of its served PDF, the filiation then
being fixed by hashes and cover pages after an external reviewer correctly objected
that the public evidence alone did not distinguish a mis-filed PDF from a duplicated
one. \S2 was found by asking, before restoring the record, whether a February claim
could re-enter an archive whose July corrections had retracted it. \textbf{Sections
3--5 came from a second external audit}, which found that the survivors list of this
note's own first draft --- Lemmas 3.11, 4.1 and A.1 --- did not survive; the same
reviewer had already caught an identical over-generous survivors list in the companion
retraction ai.viXra:2602.0084 v2. Every point was verified against the unchanged v1
proofs before adoption. That two successive drafts of two different retractions closed
their survivors lists too early is recorded here rather than quietly fixed: it is the
same failure the retractions are about.

\end{document}
